\section{Benchmark Evidence：为什么选择 tabular data，表格应该怎样读}

如果只在容易被深度网络统治的 benchmark 上测试，很难判断 dataset-level attention 是否真的有价值。作者选择 tabular data，是因为它覆盖分类与回归、连续与离散字段、数百到千万样本，同时 tree-based boosting 长期强势；这迫使 NPT 与真正有竞争力的非神经 baseline 比较。

\lecturefigure{slide-22-tabular-domain.jpg}{Tabular benchmark：规模、字段类型、任务与强基线都高度多样。}{Stanford Online 官方视频 00:37:16--00:38:07。}

\begin{knowledgebox}{基线与指标速查}
\begin{tabularx}{\textwidth}{L{0.19\textwidth}L{0.25\textwidth}>{\raggedright\arraybackslash}X}
\toprule
名称 & 解决机制 & 为什么必须比较 \\
\midrule
XGBoost / LightGBM & 高效 gradient-boosted decision trees & Tabular data 的工业强基线，善于非线性分裂与缺失/稀疏特征。 \\
CatBoost & 针对 categorical features 的 ordered boosting & 防止 NPT 只因类别字段处理更好而获得虚假优势。 \\
Random Forest / GBDT & bagging 或 boosting 的树模型 & 覆盖经典 ensemble family。 \\
MLP & 每行独立的 parametric neural baseline & 检验提升是否来自跨 datapoint mechanism。 \\
TabNet & 面向 tabular data 的 attentive neural model & 比较 domain-specific neural architecture。 \\
k-NN & 固定距离的 local lookup & 对照“learned non-parametric lookup”是否优于简单邻居。 \\
AUROC / Accuracy / RMSE & 二分类排序、多分类正确率、回归误差 & 三类任务指标方向不同，论文统一转为 method rank。 \\
\bottomrule
\end{tabularx}
\end{knowledgebox}

\teachervoice{作者选择 tabular domain 不只是因为数据方便，而是因为 general-purpose deep nets 在这里经常输给 boosting；若 NPT 只胜过一个弱 MLP，并不足以支持其通用性主张。}

\teachervoice{团队对 NPT 在 small datasets 上的 robustness 感到意外，而且没有进行极端数量的 hyperparameter 调整。这个口头观察值得记录，但它不替代严格的 tuning-budget 比较。}

\subsection{Average rank 不是 pooled accuracy}

结果页将十个 UCI datasets 按任务分组：四个 binary classification 使用 AUROC，两个 multi-class classification 使用 accuracy，四个 regression 使用 RMSE。对每个 dataset 先给 methods 排名，再在同任务组内取平均；因此数值越小越好，且不能把不同列直接当作同一物理量比较。

\lecturefigure{slide-23-tabular-benchmarks.jpg}{UCI benchmark 的平均排名与 CIFAR-10 早期结果。}{Stanford Online 官方视频 00:38:08--00:40:44。}

\begin{table}[H]
\centering
\small
\caption{论文 Table 1 的关键平均排名；数值越低越好。}
\begin{tabularx}{\textwidth}{L{0.22\textwidth}L{0.23\textwidth}L{0.23\textwidth}X}
\toprule
方法 & Binary AUROC rank & Multi-class accuracy rank & Regression RMSE rank \\
\midrule
NPT & $2.50\pm0.87$ & $2.50\pm0.50$ & $3.25\pm1.31$ \\
CatBoost & $2.75\pm0.85$ & $3.50\pm0.50$ & $3.00\pm0.91$ \\
XGBoost & $4.75\pm1.25$ & $2.50\pm1.50$ & $3.25\pm0.63$ \\
MLP & $5.75\pm1.49$ & $3.00\pm2.00$ & $5.00\pm1.22$ \\
k-NN & $8.25\pm0.48$ & $8.50\pm0.50$ & $8.75\pm0.25$ \\
\bottomrule
\end{tabularx}
\end{table}

\begin{knowledgebox}{读图：先看竞争区间，再看冠军数量}
NPT 在 binary 与 multi-class 平均 rank 位于第一梯队，在 regression 与 CatBoost/XGBoost 接近；论文还报告 NPT 在十个 tabular datasets 中四个为 top performer。更稳妥的结论是“跨多种任务具有竞争力”，而不是“统一击败所有 tree methods”。
\end{knowledgebox}

\begin{warningbox}{平均排名隐藏了什么}
十个 datasets 数量有限，不同任务组只有二到四个数据集；rank 会丢失绝对 performance gap，也受 tuning protocol 与 dataset selection 影响。标准误差来自 datasets 上的 rank variation，不是一次大样本 accuracy 的置信区间。
\end{warningbox}

\teachervoice{Q\&A 花了不少时间澄清 “4 of 10”“2 of 10” 与 standard error。讲义因此必须明确：这些数字首先描述 dataset/task 分组，其次才是 methods 的平均排名，不能从表面小数推出普遍显著优势。}

\subsection{Image result 的正确边界}

论文也在 CIFAR-10 与 MNIST 上测试 NPT。讲座展示的 CIFAR-10 结果使用 CNN encoder 加 NPT，并报告 93.7\% accuracy；这说明跨-image datapoint interaction 可以嵌入视觉 pipeline，但不能证明 NPT 取代了视觉 backbone。论文附录中的 linear patching 结果更弱，也反映 input embedding 仍会影响表现。

\begin{warningbox}{“通用输入格式”不等于“前端无关紧要”}
任何 modality 可 reshape 成 matrix，只说明接口能接收；是否保留局部结构、是否有合适 encoder、数据量与 augmentation 如何设置，仍决定 optimization 与 sample efficiency。NPT 的贡献是新增跨 datapoint mechanism，不是抹去全部 modality knowledge。
\end{warningbox}

\subsection{本章小结}

NPT 在多种 tabular tasks 上进入强基线竞争区，并展示 image extension 的可行性。证据支持“具有竞争力”，但平均 rank、有限 datasets 与不同前端不允许我们宣称它在所有场景普遍最优。

